% !TEX root = ../main.tex
%----------------------------------------------------------
% 05-native-verification.tex — Section 4: Native Possibilistic Verification
% MIGRATED from preprint A5-verification-framework.tex + A6 plain language
% Removed: appendix framing, Product 3 language
% Ozone categories → SPC categories
% Inlined plain-language explanations from A6
%----------------------------------------------------------

\section{Native Possibilistic Verification}\label{sec:native_verification}

The possibility-to-probability bridge
(Section~\ref{sec:poss_prob_bridge}) enables probabilistic verification
of possibilistic forecasts, but the conversion is lossy: it compresses
the full possibility shape into $n{+}1$ probabilities, with the entire
subnormality signal collapsed into a single ignorance bin.  Native
possibilistic verification evaluates the raw distribution directly,
preserving the shape information that the bridge discards.  The metrics
defined below complement the bridge-based lane by diagnosing forecast
quality in possibility space.

All definitions use the notation of Section~\ref{sec:poss_primer}:
$\pi$ denotes the raw (subnormal) possibility distribution,
$m = \max_{\omega \in \Omega} \pi(\omega)$ is the subnormality, and
$\ign = 1 - m$ is ignorance (Eq.~\ref{eq:ignorance}).

\begin{figure}[t]
\centering
\includegraphics[width=\columnwidth]{perf_hexbin_trajectory.png}
\caption{Possibilistic performance diagram for a three-year synthetic
    reforecast.
    \textbf{Hex bins:} colour is mean $\ign$ (dark purple = confident).
    \textbf{Contours:} joint skill $S = \alpha^* \times (1 - \eta)$.
    \textbf{Dashed diagonal:} $\delta = 0$ break-even.
    \textbf{Green line:} category-conditional means from \spc{NONE}
    (easy, upper right) to \spc{HIGH} (rare, lower left); dot
    size $\propto$ sample count; shade encodes $\ign$ (dark green
    = confident), paralleling the purple hex bins.
    \textbf{Stars:} worked scenarios from
    Section~\ref{sec:worked_examples} (green edge = hit, red = miss).
    The upper-right corner is the perfect-forecast region.}
\label{fig:performance_diagram}
\end{figure}

%----------------------------------------------------------
\subsection{Normalization Protocol}\label{sec:norm_protocol}

Some verification quantities require the raw distribution $\pi$ while
others require the normalized form $\pinorm = \pi / m$.  Two
normalization operations serve distinct purposes:

\begin{enumerate}[leftmargin=2em]
\item \textbf{Shape normalization} ($\pi \to \pinorm$): Rescales so
    that $\max(\pinorm) = 1$, isolating the system's relative
    preferences among outcomes.  Used for computing \Nc{}
    (Eq.~\ref{eq:cond_nec}) and the native scorecard
    (Section~\ref{sec:native_scorecard}).
\item \textbf{Possibility-to-probability bridge} ($\pi \to$ probability
    vector): Converts possibilities to probabilities with an explicit
    ignorance outcome (Section~\ref{sec:poss_prob_bridge}).  Used
    for information-gain verification.
\end{enumerate}

These are not alternatives: normalization operates within the possibility
space; the pignistic bridge converts to the probability space.

Table~\ref{tab:raw_norm} specifies which form each quantity uses.

\begin{table}[htbp]
\centering
\caption{Raw vs.\ normalized form used by each verification quantity.}
\label{tab:raw_norm}
\small
\renewcommand{\arraystretch}{1.15}
\begin{tabular}{@{}llcl@{}}
\toprule
Quantity & Symbol & Form & Reason \\
\midrule
\multicolumn{4}{@{}l}{\textit{Raw $(\pi)$}} \\[2pt]
\quad Ignorance              & $\ign$      & $\pi$    & Gap to normality \\
\quad Event possibility      & $\poss(A_T)$& $\pi$    & Upper bound \\
%%% U_k and L_k rows removed — IRLS appendix commented out
%%\quad Cumulative upper bound & $U_k$       & $\pi$    & For IRLS \\

\multicolumn{4}{@{}l}{\textit{Normalized $(\pinorm)$}} \\[2pt]
\quad Depth-of-truth         & $\alpha^*$  & $\pinorm$ & Shape quality \\
\quad Nonspecificity          & $\eta$      & $\pinorm$ & Shape diffuseness \\
\quad Conditional necessity   & $\Nc(A)$    & $\pinorm$ & Certainty given coverage \\[4pt]
\multicolumn{4}{@{}l}{\textit{Both}} \\[2pt]
\quad Lower threshold bound  & $L_t$       & Both     & Coverage $\times$ certainty \\
%%\quad Cumulative lower bound & $L_k$       & Both     & IRLS counterpart \\
\bottomrule
\end{tabular}
\renewcommand{\arraystretch}{1.0}
\end{table}

\textbf{Why this matters:} Getting the raw/normalized distinction wrong
is the single most common error in possibilistic scoring.  Computing
\Nc{} from raw values yields nonsense (it would be negative when
$m < 0.5$).  Computing ignorance from normalized values gives zero by
definition.
The distinction is loosely analogous to evaluating ensemble-mean error
(which uses raw ensemble members) versus a rank histogram (which uses
the normalized rank position of the observation within the ensemble).

%----------------------------------------------------------
\subsection{Five-Number Scorecard}\label{sec:native_scorecard}

Five summary statistics evaluate each forecast--observation pair, where
$c_{\mathrm{obs}}$ denotes the observed SPC category and $K = |\Omega|$
is the number of categories (here, $K = 6$).
Table~\ref{tab:scorecard_defs} lists the definitions; the equations
are given below.

\begin{table}[htbp]
\centering
\caption{Five-number scorecard definitions.}
\label{tab:scorecard_defs}
\small
\begin{tabular}{clcl}
\toprule
Symbol & Name & Component & Eq.\ \\
\midrule
$\alpha^*$ & Depth-of-truth & Possibility
    & \eqref{eq:alpha_star} \\
$\eta$ & Nonspecificity & Possibility
    & \eqref{eq:ns} \\
$\delta$ & Resolution gap & Possibility
    & \eqref{eq:rg} \\
$\ign$ & Ignorance & Ignorance
    & \eqref{eq:ignorance} \\
$\Nc^*$ & Cond.\ necessity of truth & Necessity
    & \eqref{eq:nc_star} \\
\bottomrule
\end{tabular}
\end{table}

\begin{equation}
\alpha^* = \pinorm(c_{\mathrm{obs}}) \label{eq:alpha_star}
\end{equation}
where $\alpha^* = 1$ when the truth coincides with the system's peak
category and $\alpha^* = 0$ when the system assigns zero possibility to
the observed outcome.

Following \citet{Klir1995-va}, nonspecificity measures the diffuseness
of the normalized shape:
\begin{equation}
\eta = \frac{1}{K} \sum_{c=1}^{K} \pinorm(c) \label{eq:ns}
\end{equation}
ranging from $1/K$ (sharp) to $1$ (uniform).

\begin{equation}
\delta = \alpha^* - \eta \label{eq:rg}
\end{equation}
Positive $\delta$ indicates the system assigned more possibility to the
truth than to an average category.

\begin{equation}
\ign = 1 - m = 1 - \max_{\omega \in \Omega} \pi(\omega)
\end{equation}
(same as Eq.~\ref{eq:ignorance}).

\begin{equation}
\Nc^* = \Nc(c_{\mathrm{obs}})
    = 1 - \max_{\omega \neq c_{\mathrm{obs}}} \pinorm(\omega)
\label{eq:nc_star}
\end{equation}
where $\Nc^* > 0$ only when the observed category dominates all
alternatives in the normalized distribution.  This is a direct
application of Eq.~\eqref{eq:cond_nec} evaluated at the observation.

\medskip
\noindent\textbf{\textsf{\color{rationalecolor}Plain-language glosses}}
\smallskip

\noindent
\renewcommand{\arraystretch}{1.25}
\begin{tabularx}{\columnwidth}{%
    >{\columncolor{rationalecolor!10}}c X}
$\alpha^*$ & How much normalized possibility did the system assign to
    what actually happened?  Like POD for the winning bin, but
    continuous. \\
$\eta$ & How spread out was the normalized distribution?  Like ensemble
    spread or sharpness. \\
$\delta$ & Did the system give more support to the truth than to an
    average category?  Like resolution in the Brier decomposition. \\
$\ign$ & How much did the system admit it didn't know?
    $\ign = 0$: fully confident; $\ign = 1$: complete ignorance. \\
$\Nc^*$ & Among scenarios the system covered, was the truth dominant
    enough to rule out alternatives?  A confidence flag. \\
\end{tabularx}
\renewcommand{\arraystretch}{1.0}

Aggregates $\overline{\alpha^*}$, $\overline{\eta}$,
$\overline{\delta}$, $\overline{\ign}$, $\overline{\Nc^*}$ are reported
with bootstrap confidence intervals.

Figure~\ref{fig:performance_diagram} provides a diagnostic view of the
five scorecard metrics, inspired by \citet{Roebber2009-rv}.  Hexagonal
bins aggregate forecast--observation pairs by specificity $1 - \eta$
($x$-axis) and depth-of-truth $\alpha^*$ ($y$-axis).  Bin colour
encodes mean ignorance $\ign$: dark bins are confident, pale bins
uncertain.  Hyperbolic contours mark constant joint skill
$S = \alpha^* \times (1 - \eta)$; the dashed diagonal marks
$\delta = 0$.  The green trajectory connects category-conditional
means, tracing how the system degrades from \spc{NONE} (easy, upper
right) to \spc{HIGH} (rare, lower left), with point size proportional
to sample count.

The analogy with Roebber's diagram is deliberate: specificity plays the
role of success ratio (how focused was the distribution?), $\alpha^*$
plays the role of probability of detection (did the forecast capture
truth?), and the $S$ hyperbolas play the role of CSI contours (net
quality given both axes).  The key addition is ignorance $\ign$, encoded
as bin colour, which has no deterministic analogue.

\textit{Reading the diagram:} upper right is ideal (sharp and truthful);
lower right is the worst failure mode (sharp and wrong).  Above the
$\delta = 0$ diagonal, the forecast discriminated; below it, the truth
received less support than the categorical average.  The trajectory's
slope reveals where skill degrades fastest---a steep drop between
adjacent categories signals the system's breaking point.

Both axes of Figure~\ref{fig:performance_diagram} operate on the
normalised shape $\pi'$: two forecasts with identical shapes but
different commitment levels ($\ign = 0.05$ vs.\ $\ign = 0.50$) land at
the same position.  Commitment is visible only through bin colour---a
weaker perceptual channel.
Figure~\ref{fig:commitment_diagram} promotes commitment to a primary
axis: $x = \pi_{\max} = 1 - \ign$ (how much of the unit interval the
forecaster used) versus $y = \delta = \alpha^* - \eta$ (net
discrimination).  These axes are orthogonal---$\pi_{\max}$ is the raw
\emph{scale} of $\pi$; $\delta$ is the \emph{shape quality} of $\pi'$.

\textit{Reading the commitment diagram:} upper right is committed and
discriminating (ideal); lower right is overconfident and wrong; upper
left is the cautious forecast that got the shape right.  Below the
$\delta = 0$ line, the forecast failed to discriminate.  Anchor~C
(sharp-wrong) sits squarely in the overconfident quadrant; Anchor~B
(hedged-correct) sits in the cautious-discriminating region---a
distinction invisible in the shape-quality diagram.

\FloatBarrier

\begin{figure}[t]
\centering
\includegraphics[width=\columnwidth]{commitment_diagram.png}
\caption{Commitment--discrimination diagram for the same reforecast.
    \textbf{$x$:}~commitment $\pi_{\max} = 1 - \ign$ (how much the
    forecaster staked).
    \textbf{$y$:}~discrimination $\delta = \alpha^* - \eta$ (did the
    shape favour truth?).
    These axes are orthogonal: $x$ is the raw scale of $\pi$; $y$ is
    the quality of the normalised shape $\pi'$.
    \textbf{Hex bins:} coloured by count.
    \textbf{Green line:} category means from \spc{NONE} (upper right)
    to \spc{HIGH} (lower left).
    \textbf{Stars:} Section~\ref{sec:worked_examples} anchors.
    Upper right = committed \& discriminating.
    Lower right = overconfident \& wrong.
    Dashed line = $\delta = 0$ break-even.}
\label{fig:commitment_diagram}
\end{figure}

\FloatBarrier

%----------------------------------------------------------
\subsection{Exceedance Bounds}\label{sec:exceedance_bounds}

For any threshold event $A_T \subseteq \Omega$---such as
$A_T = \{\spc{ENH},\, \spc{MDT},\, \spc{HIGH}\}$ for significant
severe weather---the possibilistic forecast produces an interval
$[L, U]$ that brackets the event's plausibility.  The bounds derive
directly from the three-component framework:

\begin{align}
U &= \poss(A_T) = \max_{\omega \in A_T} \pi(\omega)
    \label{eq:U_bound} \\
L &= m \cdot \Nc(A_T)
   = m \cdot \bigl[1 - \max_{\omega \notin A_T} \pinorm(\omega)\bigr]
    \label{eq:L_bound}
\end{align}

\noindent The interval $[L, U] = [m \cdot \Nc,\; \poss]$ answers
probability-of-exceedance questions without requiring bridge conversion:
$U$ is the maximum plausibility of any outcome in $A_T$, while $L$
combines commitment ($m$) with conditional necessity to give a lower
bound.  Width $U - L$ reflects ambiguity between ``merely possible''
and ``essentially certain.''  Under normality ($m = 1$), $L = \nec(A_T)$
and $U = \poss(A_T)$, recovering the classical necessity--possibility
interval.  The property $0 \leq L \leq U \leq 1$ always holds.

Figure~\ref{fig:upper_lower} applies these bounds to the three
worked scenarios of Section~\ref{sec:worked_examples} (Scenarios A--C),
showing each $[L, U]$ interval on a $[0,1]$ number line for
$A_T = \{\spc{ENH},\, \spc{MDT},\, \spc{HIGH}\}$.  Sharp forecasts
produce narrow intervals; hedged or ignorant forecasts produce wide
ones.

\begin{figure*}[b]
\centering
\includegraphics[width=0.75\textwidth]{upper_lower_bounds.png}
\caption{Exceedance bounds
    (Eqs.~\ref{eq:U_bound}--\ref{eq:L_bound}) for
    $A_T = \{\spc{ENH},\, \spc{MDT},\, \spc{HIGH}\}$ applied to
    Scenarios A--C (Section~\ref{sec:worked_examples}).  Each horizontal
    band shows $[L, U]$ on a $[0,1]$ number line.  Scenario~A's narrow
    interval reflects high commitment and conditional necessity;
    Scenario~B's wide interval reflects hedging; Scenario~C collapses to
    a point at zero because the system assigned zero possibility to every
    category in $A_T$.}
\label{fig:upper_lower}
\end{figure*}

\FloatBarrier
